%% ============================================================
%%  SECCIÓN 8 — Análisis Experimental
%%  Responsable: Minilux
%% ============================================================

\section{Análisis Experimental}
\label{sec:analisis_exp}

En esta sección se examina cualitativamente el comportamiento observado
en ambos métodos a partir de los datos presentados en la
Sección~\ref{sec:resultados}, sin proceder aún al cálculo de magnitudes
derivadas.

%% ------------------------------------------------------------------
\subsection{Análisis del primer método --- Tabla~\ref{tab:datos_m1}}
%% ------------------------------------------------------------------

Las cinco mediciones presentan configuraciones de jinetillos diversas
que reproducen torques de magnitud similar, lo cual es consistente con
una fuerza de tensión superficial estable bajo condiciones experimentales
equivalentes. En todos los casos se emplearon al menos dos jinetillos
distribuidos en ranuras distintas, estrategia que permite ajustar la
masa efectiva con mayor resolución que un único jinetillo, ya que la
contribución de cada uno escala con su posición relativa $n_j/10$.

Las masas efectivas se encuentran dentro del rango
$\qtyrange{2,38e-3}{2,67e-3}{\kilogram}$, con una dispersión relativa
aproximada del \qty{6}{\percent} respecto al valor central. Esta
variabilidad reducida indica buena reproducibilidad del método,
atribuible principalmente a la resolución discreta de las ranuras del
brazo mayor ($\Delta n = 1$) y a la incertidumbre en la masa de cada
jinetillo ($\Delta m = \qty{0,1}{\gram}$).

%% ------------------------------------------------------------------
\subsection{Análisis del segundo método --- Tabla~\ref{tab:datos_m2}}
%% ------------------------------------------------------------------

Los valores geométricos registrados ($a$, $b$, $h$) revelan dos aspectos
relevantes para la interpretación del resultado. En primer lugar, la
diferencia $a - b = \qty{0,18}{\centi\meter}$ es pequeña en comparación
con $h = \qty{2,78}{\centi\meter}$, de modo que el término $h^2/(a-b)$
en la Ec.~\eqref{eq:gamma_m2} alcanza aproximadamente
\qty{42,9}{\centi\meter}, muy superior a la suma
$a + b = \qty{5,28}{\centi\meter}$. En consecuencia, el resultado de
$\gamma$ queda dominado por el cociente $h^2/(a-b)$, cuya sensibilidad
frente a variaciones en $b$ es alta: un error de tan solo
$\Delta b = \qty{0,1}{\milli\meter}$ se propaga como una variación
relativa de $\sim\qty{6}{\percent}$ en $\gamma$. Esta amplificación se
ve agravada por la evaporación progresiva de la película y el drenaje
gravitacional de la solución jabonosa durante el tiempo de medición.

En segundo lugar, el uso de solución jabonosa introduce la acción de un
surfactante que se adsorbe en la interfaz líquido--aire, reduciendo las
fuerzas cohesivas entre moléculas de agua. Se espera por tanto que el
valor de $\gamma$ obtenido sea apreciablemente inferior al del primer
método y al valor bibliográfico
$\gamma_{\text{ref}} = \qty{72,8e-3}{\newton\per\meter}$
\cite{serway}, aspecto que se analiza en detalle en la
Sección~\ref{sec:conclusiones}.